% Junior Level 2, Week 7 -- workbook chapter (read aloud by a grown-up)
\chapterweek{7}{My Own Blocks}{Give a group of blocks a name and use it again.}

\section{What you will learn}
\begin{ican}
  \item I can make my own block with a name.
  \item I can use my block again and again.
  \item I can make a game easier for everyone.
\end{ican}

\section{A name for a group of blocks}
A \textbf{My Block} is a group of blocks with a name. You make it once. Then one block does all of them!

\begin{exercise}[Celebrate Block]
\begin{enumerate}[leftmargin=1.4em,itemsep=2pt]
  \item Pink-red \menu{My Blocks}: click \menu{Make a Block}.
  \item Type \enquote{celebrate}. Click \menu{OK}.
  \item Snap the blocks under \enquote{define celebrate}.
  \item Put \enquote{celebrate} under two yellow \menu{Events} blocks.
  \item Click the cat. Press space. Hooray!
\end{enumerate}
\begin{center}
\begin{scratch}[scale=0.8]
\initmoreblocks{define \namemoreblocks{celebrate}}
\blocksound{start sound \selectmenu{Pop}}
\blockrepeat{repeat \ovalnum{4}}{
  \blockmove{turn \turnright{} \ovalnum{90} degrees}
}
\blocklook{say \ovalnum{Hooray!} for \ovalnum{1} seconds}
\end{scratch}
\quad
\begin{scratch}[scale=0.8]
\blockinit{when this sprite clicked}
\blockmoreblocks{celebrate}
\end{scratch}
\quad
\begin{scratch}[scale=0.8]
\blockinit{when \selectmenu{space} key pressed}
\blockmoreblocks{celebrate}
\end{scratch}
\end{center}
\end{exercise}

\begin{exercise}[The Lost Block]
Sam's game: click the running cat to score. The flag should set the score to 0 and say \enquote{Go!}. It does not! Build it and test it.
\begin{center}
\begin{scratch}[scale=0.8]
\initmoreblocks{define \namemoreblocks{get ready}}
\blockvariable{set \selectmenu{score} to \ovalnum{0}}
\blocklook{say \ovalnum{Go!} for \ovalnum{1} seconds}
\end{scratch}
\quad
\begin{scratch}[scale=0.8]
\blockinit{when \greenflag clicked}
\blockmove{set rotation style \selectmenu{left-right}}
\blockinfloop{forever}{
  \blockmove{move \ovalnum{10} steps}
  \blockmove{if on edge, bounce}
}
\end{scratch}
\quad
\begin{scratch}[scale=0.8]
\blockinit{when this sprite clicked}
\blockvariable{change \selectmenu{score} by \ovalnum{1}}
\end{scratch}
\end{center}
Circle the block that is made but never used. Where should it go? Fix it.
\end{exercise}

\needspace{16\baselineskip}
\begin{exercise}[A Game for Everyone]
Which game is easier for everyone to play? Circle it.
\begin{center}
%% alt: Two small game screens side by side. Screen A has tiny score text and two apples that are the same shape, one red and one green. Screen B has a big score that says Score 3, a round apple to catch and a ghost to dodge.
\begin{tikzpicture}[x=1cm,y=1cm]
  \draw[line width=1pt,rounded corners=3pt] (0,0) rectangle (5.4,3.6);
  \node[font=\bfseries] at (2.7,4) {A};
  \node[font=\tiny,anchor=north west] at (0.15,3.45) {score 3};
  \fill[kcvermillion] (1.6,1.4) circle (0.45);
  \fill[kcgreen] (3.8,1.4) circle (0.45);
  \draw[line width=1pt,rounded corners=3pt] (6.4,0) rectangle (11.8,3.6);
  \node[font=\bfseries] at (9.1,4) {B};
  \node[font=\Large\bfseries,anchor=north west,fill=kcorange!25,rounded corners=2pt] at (6.55,3.45) {Score 3};
  \fill[kcvermillion] (7.9,1.2) circle (0.45);
  \node[font=\scriptsize] at (7.9,0.45) {apple};
  \node[font=\huge,kcblack] at (10.4,1.3) {\faGhost};
  \node[font=\scriptsize] at (10.4,0.45) {ghost};
\end{tikzpicture}
\end{center}
Why? Circle: \quad big words \quad shapes and colors \quad both
\end{exercise}

\begin{vocab}
  \vocabword{My Block}{A block you make yourself. Its name stands for a group of blocks.}
  \vocabword{define}{The hat block on top of a My Block. It holds the blocks your block does.}
\end{vocab}

\begin{reflection}
Draw a move with a name, like \enquote{the twirl}. How do you feel today? Circle one.

\begin{center}
\Huge\faSmile[regular]\qquad\faMeh[regular]\qquad\faFrown[regular]
\end{center}
\answerspace{3cm}
\end{reflection}
